% !TEX program = xelatex
% =============================================================================
% 中文考研复试简历 LaTeX 模板
% Author: Kody
% Repository: https://github.com/kodyyu1126/Chinese-resume-template-postgraduate
% Compiler: XeLaTeX
% =============================================================================
\documentclass[11pt,a4paper,fontset=none]{ctexart}

% ── 字体设置 ──────────────────────────────────────────────
% 若需在 Overleaf 使用自定义字体，可将已授权的宋体文件放入 fonts/ 文件夹。
% macOS 本地优先使用 Songti SC，英文字母统一使用 Times New Roman。
\IfFontExistsTF{Times New Roman}{\setmainfont{Times New Roman}}{}

\newif\ifuseuploadedcjkfonts
\IfFileExists{fonts/STSong.ttf}{
    \IfFileExists{fonts/STHeiti.ttf}{
        \IfFileExists{fonts/STKaiti.ttf}{\useuploadedcjkfontstrue}{}
    }{}
}{}

\newif\ifusesystemcjkfonts
\IfFontExistsTF{Songti SC}{\usesystemcjkfontstrue}{}

\ifuseuploadedcjkfonts
    \setCJKmainfont[
        Path = fonts/,
        BoldFont = STSong.ttf,
        ItalicFont = STSong.ttf
    ]{STSong.ttf}
    \setCJKsansfont[Path = fonts/]{STHeiti.ttf}
    \setCJKmonofont[Path = fonts/]{STHeiti.ttf}
\else\ifusesystemcjkfonts
    \setCJKmainfont[
        BoldFont = {Songti SC},
        ItalicFont = {Songti SC}
    ]{Songti SC}
    \setCJKsansfont{Songti SC}
    \setCJKmonofont{Songti SC}
\else
    \setCJKmainfont[
        BoldFont = Noto Sans CJK SC,
        ItalicFont = Noto Serif CJK SC
    ]{Noto Serif CJK SC}
    \setCJKsansfont{Noto Sans CJK SC}
    \setCJKmonofont{Noto Sans CJK SC}
\fi\fi

% ── 页面布局 ──────────────────────────────────────────────
\usepackage{geometry}
\geometry{a4paper, left=1.5cm, right=1.5cm, top=1.5cm, bottom=1.5cm}

% ── 颜色定义 ──────────────────────────────────────────────
\usepackage{xcolor}
\definecolor{sectioncolor}{RGB}{46,116,181}
\definecolor{linkcolor}{RGB}{0,112,192}
\definecolor{textcolor}{RGB}{0,0,0}

% ── 标题格式 ──────────────────────────────────────────────
\usepackage{titlesec}
\titleformat{\section}
    {\Large\bfseries\color{sectioncolor}}
    {}
    {0em}
    {\uppercase}
    [\vspace{-0.8em}\rule{\textwidth}{1pt}]
\titlespacing{\section}{0pt}{0.1em}{0.1em}

% ── 列表格式 ──────────────────────────────────────────────
\usepackage{enumitem}
\setlist[itemize]{leftmargin=*, nosep, topsep=0.2em, itemsep=0.1em, parsep=0em}

% ── 图片与排版 ────────────────────────────────────────────
\usepackage{graphicx}
\setkeys{Gin}{draft=false}
\usepackage[export]{adjustbox}
\usepackage{calc}
\usepackage{tikz}

% ── 表格 ─────────────────────────────────────────────────
\usepackage{tabularx}
\usepackage{multirow}

% ── 图标与超链接 ──────────────────────────────────────────
\usepackage{fontawesome5}
\usepackage{hyperref}
\hypersetup{colorlinks=true, linkcolor=linkcolor, urlcolor=linkcolor}

% ── 全局样式 ──────────────────────────────────────────────
\color{textcolor}
\linespread{1.15}
\pagestyle{empty}

% ── 自定义命令 ────────────────────────────────────────────
\newcommand{\daterange}[1]{{\textbf{#1}}}
\newcommand{\entrytitle}[1]{{\textbf{#1}}}

% ============================================================
\begin{document}
% ============================================================

% ── 个人信息头部 ──────────────────────────────────────────
\noindent
\begin{minipage}[t]{\textwidth}
    \vspace*{-2.5em}
    \begin{center}
        {\Huge\bfseries\color{sectioncolor} 葛飞翔} \\[0.9em]
        \textbf{电话：}19857950790\hspace{1.5em}
        \textbf{邮箱：}\href{mailto:marshall_gefxiang@163.com}{marshall\_gefxiang@163.com} \\[0.45em]
        \textbf{研究兴趣：}大模型分布式系统；深度学习算法
    \end{center}
\end{minipage}

% 头像叠加在页面右上角，不影响顶部信息的居中对齐。
\begin{tikzpicture}[remember picture, overlay]
    \node[anchor=north east, inner sep=0pt, xshift=-1.5cm, yshift=-1.5cm]
        at (current page.north east) {\includegraphics[width=2.8cm,height=3.3cm,keepaspectratio]{assets/profile-photo.png}};
\end{tikzpicture}

\vspace{0.3em}

% ── 教育经历 ──────────────────────────────────────────────
\section*{教育经历}
\vspace{-0.3em}

\noindent
\entrytitle{中国科学院大学｜推免硕士研究生}
\hfill\daterange{2025.09 \textasciitilde{} 至今}

\par\noindent
\textbf{人工智能}

\noindent
\entrytitle{上海对外经贸大学｜本科}
\hfill\daterange{2021.09 \textasciitilde{} 2025.06}

\par\noindent
\textbf{数据科学与大数据技术}

\par\noindent
$\bullet$ \textbf{学业表现：}GPA：3.74/4.00；专业排名：1/39；综测排名：2/153；CET-4：615；CET-6：554。

\par\noindent
$\bullet$ \textbf{相关课程：}算法导论（97）、人工智能（96）、深度学习（96）、数字图像处理（95）、分布式计算（95）。

% ── 研究工作 ──────────────────────────────────────────────
\section*{研究工作}
\vspace{-0.1em}

\noindent
\entrytitle{Using Domain Adapter and Prompt Learner for Few-shot Teaching Action Recognition in videos}

\par\noindent
$\bullet$ \textbf{论文信息：}Feixiang Ge, Xinlei Li, Yan Rong, et al.，IJCNN 2025，Accepted, to be published；第一作者，CCF-C。

\par\noindent
$\bullet$ \textbf{方法贡献：}提出基于 CLIP 视觉-语言大模型的 TBA-CLIPNet，面向小样本及零样本识别；引入领域适配器弥补预训练大模型与下游任务间的语义鸿沟，构建提示模板学习器从图像特征中自动学习提示模板；设计教师目标识别算法 TD，并利用时空自注意力机制提取视频时序信息。

\par\noindent
\entrytitle{A Hybrid CNN-Transformer Network for Fine-Grained Tongue Multi-Attribute Classification with Tongue Group Attribute Mask Training in Traditional Chinese Medicine Diagnosis}

\par\noindent
$\bullet$ \textbf{论文信息：}Liu J, Ge F, Li F, et al.，International Conference on Intelligent Computing，Springer Nature Singapore，2025：467--479；共同一作，CCF-C。

\par\noindent
\entrytitle{可拓基元提取以及可拓知识论坛系统}
\hfill\daterange{软件著作权}

\par\noindent
$\bullet$ \textbf{登记号：}2024SR0643195；第一著作权人。

% ── 业界工作 ──────────────────────────────────────────────
\section*{业界工作}
\vspace{-0.1em}

\noindent
\entrytitle{百度｜混合云模型适配组｜模型适配与推理性能优化}
\hfill\daterange{2026.02 \textasciitilde{} 至今}

\par\noindent
\entrytitle{DeepSeek-V3.2 单机与 PD 分离推理调优}

\par\noindent
$\bullet$ \textbf{性能优化：}面向昆仑芯 P800，基于 SGLang 与 vLLM-Kunlun 对齐算子链路，围绕 NSA Indexer、MTP/投机解码和上下文并行开展细粒度打点、框架逻辑对齐与关键函数优化；相关 NSA Indexer BF16 优化已合入。

\par\noindent
$\bullet$ \textbf{量化结果：}单机性能较原准出版本大幅提升，在不同并发与输入长度场景均取得显著收益。PD 分离 Router 方案较上一版提升 10\%--70\%，上下文并行使短文本性能平均提升约 30\%；定位投机解码导致的乱码问题并通过关闭相关路径恢复精度。

\par\noindent
\entrytitle{P800 大模型适配与客户场景准出}

\par\noindent
$\bullet$ \textbf{模型适配：}完成 Qwen3/Qwen3.5 系列模型的推理适配、性能精度验证和准出文档；定位 prompt\_logprobs 算子兼容、vLLM 版本差异、prefix cache、CUDA Graph、xgrammar 与 XCCL 等问题。

\par\noindent
$\bullet$ \textbf{场景覆盖：}参与 ASR、Grounding SAM、OpenFold/SAM3、BEVFormer 等模型在 P800 上的训练或推理适配；针对长文本与超长上下文需求，推进显存管理、量化模型和推理链路优化。

\par\noindent
\entrytitle{vLLM-Omni 升级与多模态推理优化}

\par\noindent
$\bullet$ \textbf{框架升级：}重构 Kunlun Plugin 与开源 vLLM-Omni 的 patch 解耦方式，同步上游特性，输出标准镜像与 SOP；完成 Wan、Qwen-Image、GLM-Image、Qwen-Image-Edit 等多个模型的适配和初步精度/性能准出，整体调优提升 50\%--100\%+。

\par\noindent
$\bullet$ \textbf{性能结果：}Wan 系列文本到视频模型推理耗时大幅下降；GLM-Image 达到 H20 的 80\%--100\% 性能且精度符合预期，Qwen-Image-Edit 达到 H20 的 50\%。

\par\noindent
\newpage
\entrytitle{Dynamo 分布式推理与 HiCache 性能分析}

\par\noindent
$\bullet$ \textbf{工程落地：}完成 Dynamo 源码编译与部署，搭建 SGLang + HiCache + Dynamo 环境，验证 GLM 系列模型；修复 KV-aware 模式 DP rank 不一致导致的 Hang，解决前缀复用调度失效问题：修正 SGLang KV event 发布逻辑，完善 Dynamo 侧 event key（worker\_id、DP rank）管理；完成 HiCache 端到端测试，并调研 fast tokenization。

\par\noindent
$\bullet$ \textbf{性能分析：}在 Dynamo 与 Router 两侧注入 Prometheus 指标，拆解 tokenize、transport、dispatch、route 阶段；GLM 系列测试显示 TPOT 稳定提升 10\%+，短文本 TTFT 提升 5\%--20\%，总吞吐有收益，并形成 Dynamo 方案进度与准出材料。

\par\noindent
\entrytitle{vLLM-Kunlun 开源贡献与平台提效}

\par\noindent
$\bullet$ \textbf{开源贡献：}围绕 vLLM-Kunlun Workspace Manager 与昆仑芯显存管理提交并合入 \href{https://github.com/baidu/vLLM-Kunlun/pull/283/changes}{PR 283}；协助完成量化 patch 对齐及 Triton 算子替换（\href{https://github.com/baidu/vLLM-Kunlun/pull/409}{PR 409}、\href{https://github.com/baidu/vLLM-Kunlun/pull/351}{PR 351}、\href{https://github.com/baidu/vLLM-Kunlun/pull/381}{PR 381}）。

\par\noindent
$\bullet$ \textbf{标准化与赋能：}迭代 LLM 部署 SOP Skill 与准出 Agent，自动化模型部署、性能/精度测试和结构化报告生成；编写 vLLM-Omni、DiT 优化、Dynamo 进度等技术材料，提升模型适配与交付效率。

\noindent
\entrytitle{eBay 软件工程有限公司｜AI Platform｜机器学习工程师（Intern）}
\hfill\daterange{2024.10 \textasciitilde{} 2025.08}

\par\noindent
$\bullet$ \textbf{平台工程：}深入实践 Ray Data、Ray Serve 等组件，解决多类型资源标签调度和 PyArrow 内存泄漏问题；相关修复已合入公司内部 Ray 分支（\href{https://github.com/ray-project/ray/pull/49124}{PR 49124}、\href{https://github.com/ray-project/ray/pull/49756}{PR 49756}）。

\par\noindent
$\bullet$ \textbf{推理与工作流：}修复 KubeRay 资源调度缺陷使 GPU 利用率翻倍，参与 AI 工作流框架开发，并探索 LLM 强化学习训练的资源需求、GPU 利用率与数据吞吐。

% ── 荣誉情况 ──────────────────────────────────────────────
\section*{荣誉情况}
\vspace{-0.1em}

\noindent
\begin{tabularx}{\textwidth}{@{}X X@{}}
    $\bullet$ 2024 年度上海市奖学金 & $\bullet$ 第十四届蓝桥杯国赛三等奖（国家级） \\
    $\bullet$ 全国大学生市场调查分析大赛三等奖（国家级） & $\bullet$ 上海对外经贸大学优秀毕业生 \\
    $\bullet$ 美国大学生数学建模竞赛 H 奖（国际级） & $\bullet$ 2023--2024 爱建特种基金会奖学金 \\
    $\bullet$ 国科大杭州高等研究院新生奖学金 & $\bullet$ 上海对外经贸大学优秀共青团员
\end{tabularx}

\end{document}
